\chapter{Fertility Awareness}

\section*{Definition and Overview}
Fertility awareness is the practice of identifying the fertile and infertile phases of the menstrual cycle by observing the body's natural signs, and using that knowledge either to avoid pregnancy (as a method of contraception) or to plan a pregnancy. The methods are based on tracking cervical mucus, basal body temperature, and cycle calendar patterns, sometimes combined as the symptothermal method. Fertility awareness requires commitment, daily observation, and a fairly regular cycle to be reliable, and its effectiveness as contraception varies widely with how strictly it is used. It is free of hormones and side effects, which appeals to many women, and it is also a valuable tool for couples trying to conceive. This chapter explains how the methods work, their reliability, and how to use them properly.

\section*{Epidemiology}
Fertility awareness is used by a minority of couples as their primary contraceptive method, but interest has grown in recent years with the rise of cycle-tracking apps. As a family planning method, its effectiveness depends heavily on the user: with perfect, strict use, typical failure rates of around 2--4\% per year are achievable, but with typical use, failure rates of around 10--20\% per year are reported. Cycle-tracking apps vary widely in quality and accuracy, and several popular apps have been shown to be inaccurate in predicting ovulation. As a fertility tool, awareness of the fertile window is central to natural family planning and to timing intercourse for conception.

\section*{Aetiology and Pathophysiology}
Fertility awareness works by detecting the physiological changes that mark the fertile window.
\begin{itemize}
    \item \textbf{The fertile window:} a woman is fertile for roughly 6 days each cycle, the 5 days before ovulation plus the day of ovulation itself, because sperm can survive for up to 5 days
    \item \textbf{Cervical mucus:} rising oestrogen before ovulation makes mucus increase, become clearer, stretchy, and slippery (like raw egg white), signalling fertility; after ovulation, progesterone makes mucus thick, sticky, and cloudy
    \item \textbf{Basal body temperature:} progesterone released after ovulation raises the resting body temperature by about 0.2--0.5 degrees, which stays elevated until the next period; the rise confirms that ovulation has occurred
    \item \textbf{Calendar calculation:} the cycle is charted to estimate the fertile window, but this alone is the least reliable sign because cycle length varies
    \item \textbf{Cervical position:} some methods also track the position and firmness of the cervix, which changes around ovulation
    \item \textbf{Ovulation timing:} ovulation usually occurs about 14 days before the next period, but its exact timing varies between cycles
    \item \textbf{Luteinising hormone (LH) tests:} home ovulation kits detect the LH surge 24--36 hours before ovulation, useful for conception
\end{itemize}

\section*{Clinical Features}
The features are the observed daily signs.
\begin{itemize}
    \item Changes in cervical mucus: from dry to sticky, then wet, slippery, and stretchy around ovulation, then sticky again
    \item A rise in basal body temperature, measured first thing each morning before getting out of bed
    \item Mid-cycle pain or twinges in some women around ovulation
    \item Breast tenderness and other premenstrual changes in the luteal phase
    \item The fertile window marked on a chart or in an app
    \item For contraception: abstinence or barrier use during the fertile window
    \item For conception: intercourse timed to the fertile days
\end{itemize}

\section*{Differential Diagnosis}
Considerations when choosing or using fertility awareness include:
\begin{itemize}
    \item Cycle irregularity, which makes the methods unreliable; cycles are commonly irregular in adolescence, perimenopause, and conditions such as polycystic ovary syndrome
    \item Confounding factors that affect the signs: fever and illness raise temperature; some medicines, including the pill, abolish the natural signs; breastfeeding and recent childbirth disrupt cycles
    \item Vaginal discharge from infection, which can be confused with fertile mucus
    \item The need to rule out pregnancy when a period is late
    \item Alternative contraceptives, including long-acting reversible methods that are far more forgiving of human error
\end{itemize}

\section*{When to Seek Medical Care}
Women considering fertility awareness for contraception should discuss it with a doctor or family planning service, especially if cycles are irregular, if they are in a life stage where a pregnancy would be a problem, or if they need help learning the methods. Couples trying to conceive should seek advice if they are not pregnant after 12 months of trying (6 months if over 35).

\textbf{Call 000 (or local emergency services) for:}
\begin{itemize}
    \item Severe abdominal or pelvic pain
    \item Vaginal bleeding with pain and faintness in early pregnancy, which may indicate an ectopic pregnancy
    \item Any emergency symptoms during pregnancy
\end{itemize}

\section*{Diagnosis and Investigations}
\begin{itemize}
    \item \textbf{Teaching and charting:} learning the methods from a trained educator, with daily charting of mucus, temperature, and cycle days
    \item \textbf{Symptothermal method:} the most reliable form, combining at least two signs (usually mucus and temperature)
    \item \textbf{Ultrasound:} used in fertility assessment to track follicle development and confirm ovulation
    \item \textbf{Progesterone blood test:} can confirm ovulation in mid-cycle
    \item \textbf{LH testing:} for timing intercourse or insemination
    \item \textbf{App selection:} choose apps based on validated methodology; many popular apps are not accurate enough for contraception
\end{itemize}

\section*{Management}
\begin{itemize}
    \item \textbf{For contraception:} chart daily, identify the fertile window, and use abstinence or a barrier method during it; combine methods (symptothermal) for reliability
    \item \textbf{For conception:} time intercourse to the fertile days, and use LH kits or an app alongside the signs if desired
    \item \textbf{During irregular cycles:} rely more heavily on mucus and temperature signs rather than calendar predictions, or choose a different method
    \item \textbf{Back-up contraception:} many couples keep condoms or another method for the fertile window
    \item \textbf{Emergency contraception:} know how to access it quickly if unprotected sex occurs in the fertile window
    \item \textbf{Combined with other methods:} fertility awareness is compatible with barrier methods and with fertility treatments
\end{itemize}

\section*{Complications}
\begin{itemize}
    \item Unintended pregnancy from method failure, especially with calendar-only methods or inconsistent charting
    \item Delays in conception from mistimed intercourse in couples using the method for family planning
    \item Anxiety and reduced spontaneity in sexual relationships
    \item Reliance on inaccurate apps leading to pregnancy or to false reassurance
    \item No protection against sexually transmitted infections
\end{itemize}

\section*{Prognosis and Outcomes}
For couples who are committed to daily charting and have regular cycles, fertility awareness can be an effective method of contraception, and it is among the few that involve no hormones or devices. Its reliability depends on the user and the method used, and it is considerably less forgiving than long-acting reversible contraception. As a fertility tool, understanding the fertile window significantly improves the odds of conception when intercourse is timed correctly. The best outcomes come from good teaching, honest assessment of one's own reliability, and a willingness to use back-up methods.

\section*{Prevention and Self-Care}
\begin{itemize}
    \item Learn the methods from a trained educator or accredited app rather than relying on guesswork
    \item Chart mucus and temperature daily at the same times, and record every day
    \item Treat any fever as a reason to distrust temperature readings that day
    \item Avoid relying on calendar predictions alone, especially with irregular cycles
    \item Keep emergency contraception knowledge and supplies in mind
    \item If trying to conceive, seek help if pregnancy has not occurred within the recommended timeframe
    \item Use condoms in addition if you need protection from sexually transmitted infections
\end{itemize}

\section*{Sources and Further Reading}
The following reputable sources provide further information for healthcare students and patients seeking reliable, current advice about fertility awareness.
\begin{itemize}
    \item Family Planning NSW. \textit{Fertility awareness fact sheet.} https://www.fpnsw.org.au/
    \item Healthdirect Australia. \textit{Fertility awareness methods.} https://www.healthdirect.gov.au/
    \item Your Fertility. \textit{Understanding your fertile window.} https://www.yourfertility.org.au/
    \item NHS. \textit{Fertility awareness methods.} https://www.nhs.uk/conditions/contraception/
    \item Mayo Clinic. \textit{Fertility awareness.} https://www.mayoclinic.org/
    \item World Health Organization. \textit{Family planning.} https://www.who.int/
\end{itemize}
